\section{Introduction}
\label{sec:introduction}

Remaining Useful Life (RUL) prediction estimates how many more operating
cycles a machine can sustain before failure.  Accurate RUL estimates
enable condition-based maintenance, reducing both unplanned failures and
unnecessary preventive replacements~\cite{saxena_2008_damage_prognosis}.

We study the NASA Commercial Modular Aero-Propulsion System Simulation
(C-MAPSS) FD001 dataset, a widely used benchmark comprising 100 training
and 100 test turbofan engines.  Each engine is monitored by 21 sensors
and 3 operational settings over its full run-to-failure trajectory.
After removing 7 constant or near-constant sensors, 14 informative
channels remain.

\subsection{Phase 1 Recap}
In Phase~1 we constructed hand-crafted rolling-window features (mean,
standard deviation and slope over the most recent 10 cycles) for each
sensor channel, yielding 57 features per observation.  Four regression
models were trained on 80 engines and evaluated on 100 fixed validation
points from 20 held-out engines:

\begin{table}[h]
\centering
\caption{Phase 1 validation results (100 fixed points)}
\label{tab:phase1_results}
\begin{tabular}{lcc}
\toprule
Model & MAE & RMSE \\
\midrule
OLS & 32.86 & 42.92 \\
Elastic Net & 31.66 & 41.77 \\
Random Forest & 25.03 & 40.17 \\
Gradient Boosting & 26.12 & 40.17 \\
\bottomrule
\end{tabular}
\end{table}

Key weaknesses identified in Phase~1: (a) systematic under-prediction
for engines with true RUL above 200 cycles (mean error approximately
$-120$ cycles for linear models), (b) the 10-cycle rolling window
captures only very recent sensor behaviour, and (c) RMSE substantially
exceeds MAE, indicating a few large errors dominate.

\subsection{Phase 2 Contributions}
Phase~2 addresses these weaknesses with four contributions:
\begin{enumerate}
  \item \textbf{Sequence models.} LSTM and GRU networks learn temporal
    patterns directly from 30-cycle sliding windows, removing the need
    for hand-crafted features.
  \item \textbf{Alternative architectures.} A 1D-CNN and a Transformer
    encoder provide complementary inductive biases (local convolution
    versus global self-attention).
  \item \textbf{Hybrid stacking.} A second-stage Random Forest combines
    deep-learning predictions with traditional rolling features in an
    out-of-fold stacking framework.
  \item \textbf{Ablation studies.} Seven controlled experiments quantify
    the impact of window length, input features, RUL capping, model
    capacity, regularisation, architecture choice and hybrid components.
\end{enumerate}

The rest of this paper is organised as follows.
Section~\ref{sec:literature} reviews related work.
Section~\ref{sec:dataset} describes the dataset and preprocessing.
Sections~\ref{sec:recurrent}--\ref{sec:hybrid} present the four model
families.  Section~\ref{sec:ablations} reports the ablation studies.
Section~\ref{sec:results} gives the final evaluation on both validation
and held-out test sets with bootstrap confidence intervals.
Section~\ref{sec:conclusion} concludes.
